%%%%%%%%%%%%%%%%%%%%%%%%%%%%%%%%%%%%%%%%%%%%%%%%%%%%%%%%%%%%%%%%%%%%%%%%%%%%%%%
% APPENDIX U: Kahneman & Tversky – Systematic Framework Integration
% The Psychological Microfoundations of Context-Dependent Choice
%%%%%%%%%%%%%%%%%%%%%%%%%%%%%%%%%%%%%%%%%%%%%%%%%%%%%%%%%%%%%%%%%%%%%%%%%%%%%%%
\section{Kahneman \& Tversky: The Psychological Microfoundations}
\label{app:kahneman-tversky}

%==============================================================================
% PART 0: INTEGRATION OVERVIEW
%==============================================================================
\subsection{Framework Integration Map}

This appendix demonstrates how the research program of Daniel Kahneman and 
Amos Tversky provides the \textbf{psychological microfoundations} for the 
$(C, \Psi, C^*, K, Q)$ framework. Their work establishes why context ($\Psi$) 
must enter utility functions and why standard expected utility theory fails 
as a descriptive model.

\begin{center}
\renewcommand{\arraystretch}{1.4}
\begin{tabular}{|l|l|l|}
\hline
\textbf{Framework Element} & \textbf{K/T Contribution} & \textbf{Key Papers} \\
\hline
$U = U(x - r)$ (Reference Dependence) & Utility is relative, not absolute & PT 1979, Tversky \& Kahneman 1991 \\
$\lambda > 1$ (Loss Aversion) & Asymmetric value function & PT 1979, Kahneman et al. 1990 \\
$\Psi_{frame}$ (Framing Effects) & Context determines choice & Tversky \& Kahneman 1981 \\
$\Psi_C$ (Cognitive Dimension) & Bounded rationality operationalized & Kahneman 2011, TJ Heuristics 1974 \\
$C^*$ Convergence Mechanism & System 1 vs. System 2 & Kahneman 2003, 2011 \\
\hline
\end{tabular}
\end{center}

\paragraph{Kahneman/Tversky's Unique Role.}
While other contributors address specific framework components (Fehr: $U^{SOC}$; 
Bloom: $\Psi_T$; Acemoglu: $\Psi_I$), Kahneman and Tversky uniquely provide:
\begin{enumerate}
    \item \textbf{The fundamental critique} of expected utility as descriptive theory
    \item \textbf{Reference dependence} as the core mechanism for context-sensitivity
    \item \textbf{Loss aversion} as the empirically dominant asymmetry
    \item \textbf{Dual-process theory} explaining when $C^*$ is reached vs. violated
\end{enumerate}

%==============================================================================
% PART I: REFERENCE DEPENDENCE – U(x-r) INSTEAD OF U(x)
%==============================================================================
\subsection{Part I: Reference Dependence – The Core Innovation}

\subsubsection{Theoretical Foundation}

Standard expected utility theory posits:
\begin{equation}
U = U(x) \quad \text{where } x = \text{final wealth state}
\end{equation}

Kahneman and Tversky's fundamental insight: utility is defined over \textit{changes} 
from a reference point, not absolute levels:
\begin{equation}
\boxed{
U = v(x - r) \quad \text{where } r = r(\Psi) \text{ is context-dependent}
}
\label{eq:kt-reference}
\end{equation}

\paragraph{Framework Implication.}
The reference point $r$ is itself a function of context $\Psi$:
\begin{itemize}
    \item Status quo ($r = x_0$)
    \item Expectations ($r = \mathbb{E}[x]$)
    \item Social comparison ($r = x_{peers}$)
    \item Aspiration levels ($r = x_{aspiration}$)
\end{itemize}

This means: \textbf{The same outcome $x$ produces different utility depending on $\Psi$}.

\subsubsection{The Prospect Theory Value Function}

From Kahneman \& Tversky (1979):
\begin{equation}
v(x - r) = 
\begin{cases}
(x - r)^\alpha & \text{if } x \geq r \\
-\lambda (r - x)^\beta & \text{if } x < r
\end{cases}
\label{eq:kt-value}
\end{equation}

With empirically estimated parameters:
\begin{itemize}
    \item $\alpha \approx \beta \approx 0.88$ (diminishing sensitivity)
    \item $\lambda \approx 2.25$ (loss aversion coefficient)
\end{itemize}

\subsubsection{Framework Translation}

Reference dependence directly implies context-dependent utility:
\begin{equation}
\frac{\partial U}{\partial \Psi} \neq 0 \quad \text{via} \quad \frac{\partial r}{\partial \Psi} \neq 0
\end{equation}

This is the \textbf{microfoundation} for why $\Psi$ enters all framework equations.

%==============================================================================
% PART II: LOSS AVERSION – λ > 1 IN FEPSDE
%==============================================================================
\subsection{Part II: Loss Aversion – The Dominant Asymmetry}

\subsubsection{Theoretical Foundation}

Loss aversion states that losses loom larger than equivalent gains:
\begin{equation}
\boxed{
|v(-\Delta)| > |v(+\Delta)| \quad \Leftrightarrow \quad \lambda > 1
}
\label{eq:kt-loss-aversion}
\end{equation}

\paragraph{Empirical Magnitude.}
Meta-analyses consistently find $\lambda \in [1.5, 2.5]$, with $\lambda \approx 2$ 
as central estimate. This means: \textbf{A loss hurts roughly twice as much as 
an equivalent gain pleases.}

\subsubsection{Integration into FEPSDE Structure}

The framework's FEPSDE welfare function incorporates loss aversion:
\begin{equation}
U = \sum_{d \in \{F,E,P,S,D,E\}} \delta_d^t \left[ G_{d,t} - \lambda_d \cdot L_{d,t} \right]
\label{eq:fepsde-loss}
\end{equation}

where:
\begin{itemize}
    \item $G_{d,t}$ = gains in dimension $d$ at time $t$
    \item $L_{d,t}$ = losses in dimension $d$ at time $t$
    \item $\lambda_d > 1$ = dimension-specific loss aversion
\end{itemize}

\paragraph{Empirical Evidence by Domain.}

\begin{center}
\renewcommand{\arraystretch}{1.3}
\begin{tabular}{|l|c|l|}
\hline
\textbf{FEPSDE Dimension} & \textbf{$\lambda$ Estimate} & \textbf{Source} \\
\hline
Financial ($F$) & $2.0 - 2.5$ & Kahneman et al. 1990, Benartzi \& Thaler 1995 \\
Emotional ($E$) & $1.5 - 2.0$ & Baumeister et al. 2001 ("Bad is stronger than good") \\
Physical ($P$) & $2.0 - 3.0$ & Chapman 1996 (health states) \\
Social ($S$) & $2.5 - 3.5$ & Baumeister \& Leary 1995 (social exclusion) \\
Digital ($D$) & $1.5 - 2.0$ & Acquisti et al. 2013 (privacy) \\
Ecological ($E$) & $2.0 - 2.5$ & Kahneman \& Knetsch 1992 (environmental goods) \\
\hline
\end{tabular}
\end{center}

\subsubsection{Policy Implication}

Loss aversion explains why:
\begin{itemize}
    \item Status quo bias is pervasive ($C^*$ is sticky)
    \item Reforms framed as "preventing losses" outperform "achieving gains"
    \item Endowment effects create apparent $C_{ij}$ where none exists objectively
\end{itemize}

%==============================================================================
% PART III: FRAMING EFFECTS – Ψ_frame AS CONTEXT DIMENSION
%==============================================================================
\subsection{Part III: Framing Effects – $\Psi_{frame}$ as Context}

\subsubsection{Theoretical Foundation}

Framing effects demonstrate that \textit{logically equivalent descriptions} 
produce different choices:
\begin{equation}
\boxed{
\text{Choice}(A, \Psi_{frame}^1) \neq \text{Choice}(A, \Psi_{frame}^2) 
\quad \text{even when } A^1 \equiv A^2
}
\label{eq:kt-framing}
\end{equation}

\subsubsection{The Classic Asian Disease Problem}

Tversky \& Kahneman (1981):

\textbf{Positive Frame:}
\begin{quote}
Program A: 200 people will be saved. \\
Program B: 1/3 probability 600 saved, 2/3 probability nobody saved.
\end{quote}
Result: 72\% choose A (risk-averse in gains domain)

\textbf{Negative Frame:}
\begin{quote}
Program C: 400 people will die. \\
Program D: 1/3 probability nobody dies, 2/3 probability 600 die.
\end{quote}
Result: 78\% choose D (risk-seeking in losses domain)

\paragraph{Framework Translation.}
Programs A=C and B=D are identical. The reversal demonstrates:
\begin{equation}
\Psi_{frame} \in \Psi \quad \text{and} \quad \frac{\partial \text{Choice}}{\partial \Psi_{frame}} \neq 0
\end{equation}

\subsubsection{Framing as $\Psi$ Subcomponent}

Framing effects operate through reference point manipulation:
\begin{equation}
\Psi_{frame} \to r(\Psi_{frame}) \to v(x - r) \to \text{Choice}
\end{equation}

This suggests framing is not a separate dimension but operates \textit{through} 
the existing $\Psi$ structure—primarily via $\Psi_K$ (information/presentation) 
and $\Psi_C$ (cognitive processing).

%==============================================================================
% PART IV: HEURISTICS AND BIASES – Ψ_C OPERATIONALIZATION
%==============================================================================
\subsection{Part IV: Heuristics and Biases – $\Psi_C$ Operationalization}

\subsubsection{Theoretical Foundation}

The Heuristics and Biases program (Tversky \& Kahneman 1974) establishes that 
cognitive processing ($\Psi_C$) systematically deviates from Bayesian rationality:

\begin{equation}
\boxed{
\text{Judgment} = f(\text{Information}, \Psi_C) \neq \text{Bayes}(\text{Information})
}
\label{eq:kt-heuristics}
\end{equation}

\subsubsection{Key Heuristics and Framework Translation}

\begin{center}
\renewcommand{\arraystretch}{1.4}
\small
\begin{tabular}{|p{2.5cm}|p{4cm}|p{5cm}|}
\hline
\textbf{Heuristic} & \textbf{Definition} & \textbf{Framework Translation} \\
\hline
Availability & Probability judged by ease of recall & $\Psi_K$ (information salience) affects perceived $C_{ij}$ \\
\hline
Representativeness & Probability by similarity to prototype & $\Psi_C$ distorts updating; base rates ignored \\
\hline
Anchoring & Estimates biased toward initial values & Reference point $r$ is path-dependent \\
\hline
Affect Heuristic & Risk judged by emotional response & $U^E$ contaminates probability estimates \\
\hline
\end{tabular}
\end{center}

\subsubsection{Two Systems Theory}

Kahneman (2011) systematizes these findings:

\begin{center}
\begin{tabular}{|l|l|l|}
\hline
& \textbf{System 1} & \textbf{System 2} \\
\hline
Operation & Automatic, fast & Controlled, slow \\
Effort & Low & High \\
Awareness & Unconscious & Conscious \\
Errors & Systematic biases & Occasional lapses \\
\hline
\textbf{Framework} & Default $C$ processing & Deliberate $C^*$ seeking \\
\hline
\end{tabular}
\end{center}

\paragraph{Framework Implication.}
Convergence to $C^*$ depends on System 2 engagement:
\begin{equation}
K \to K^* \quad \text{iff} \quad \Psi_C \text{ supports System 2 activation}
\end{equation}

This explains heterogeneity in coherence: high-$\Psi_C$ contexts (time, motivation, 
cognitive resources) enable $C^*$ convergence; low-$\Psi_C$ contexts produce 
systematic deviations.

%==============================================================================
% PART V: SYNTHESIS – THE KAHNEMAN/TVERSKY EQUATIONS
%==============================================================================
\subsection{Part V: Synthesis – The Kahneman/Tversky Equations}

\subsubsection{Complete Framework Translation}

\begin{center}
\renewcommand{\arraystretch}{1.5}
\begin{tabular}{|l|c|l|}
\hline
\textbf{K/T Finding} & \textbf{Equation} & \textbf{Framework Element} \\
\hline
Reference dependence & $U = v(x - r(\Psi))$ & Context enters utility directly \\
\hline
Loss aversion & $\lambda \approx 2 > 1$ & FEPSDE asymmetry parameter \\
\hline
Diminishing sensitivity & $\alpha, \beta < 1$ & Concavity in gains, convexity in losses \\
\hline
Probability weighting & $\pi(p) \neq p$ & Overweight small $p$, underweight large $p$ \\
\hline
Framing effects & $\partial \text{Choice}/\partial \Psi_{frame} \neq 0$ & $\Psi_K$ affects decisions \\
\hline
Availability & $P_{judged} = g(\text{salience})$ & $\Psi_K$ distorts perceived $C_{ij}$ \\
\hline
Anchoring & $r = f(r_0, \text{adjustment})$ & Reference point is path-dependent \\
\hline
System 1 / System 2 & $K \to K^*$ requires System 2 & $\Psi_C$ moderates coherence \\
\hline
\end{tabular}
\end{center}

\subsubsection{Kahneman/Tversky's Unique Contribution}

\textbf{Without Kahneman/Tversky, the framework would:}
\begin{itemize}
    \item Lack the psychological \textit{mechanism} for why $\Psi$ matters
    \item Have no microfoundation for reference-dependent utility
    \item Miss the $\lambda > 1$ asymmetry in FEPSDE
    \item Cannot explain why identical options produce different choices (framing)
    \item Lack a theory of \textit{when} agents converge to $C^*$ vs. deviate
\end{itemize}

\textbf{With Kahneman/Tversky, the framework gains:}
\begin{itemize}
    \item Psychological realism: utility is relative, not absolute
    \item Quantified parameters: $\lambda \approx 2$, $\alpha \approx 0.88$
    \item Mechanism for $\Psi$-dependence: reference points shift with context
    \item Dual-process theory: $\Psi_C$ determines $C^*$ convergence
    \item Empirical toolkit: decades of experimental methodology
\end{itemize}

%==============================================================================
% PART VI: COMPLETE PAPER DOCUMENTATION
%==============================================================================
\subsection{Part VI: Complete Paper Documentation}

\subsubsection{Foundational Papers}

\paragraph{Paper 1: Kahneman \& Tversky (1979)}

``Prospect Theory: An Analysis of Decision under Risk.''
\textit{Econometrica}, 47(2), 263--291.

\textbf{Summary.}
The foundational paper establishing Prospect Theory. Demonstrates systematic 
violations of expected utility theory: certainty effect, reflection effect, 
isolation effect. Introduces the S-shaped value function with reference dependence 
and loss aversion.

\textbf{Framework Translation.}
\begin{equation}
U_{EU} = \sum_i p_i u(x_i) \quad \to \quad U_{PT} = \sum_i \pi(p_i) v(x_i - r)
\end{equation}
This is the core equation showing how $\Psi$ (via $r$) enters utility.

\textbf{Key Finding.}
Most cited economics paper of all time. Establishes that $U \neq U(x)$ but 
$U = U(x - r)$ where $r$ depends on context.

\paragraph{Paper 2: Tversky \& Kahneman (1974)}

``Judgment under Uncertainty: Heuristics and Biases.''
\textit{Science}, 185(4157), 1124--1131.

\textbf{Summary.}
Documents three fundamental heuristics (representativeness, availability, 
anchoring) and their associated biases. Establishes that human judgment 
systematically deviates from Bayesian norms.

\textbf{Framework Translation.}
Provides the empirical content for $\Psi_C$: cognitive limitations are 
systematic and predictable, not random noise.

\paragraph{Paper 3: Tversky \& Kahneman (1981)}

``The Framing of Decisions and the Psychology of Choice.''
\textit{Science}, 211(4481), 453--458.

\textbf{Summary.}
Demonstrates that logically equivalent descriptions produce different choices.
The Asian Disease problem shows risk preferences reverse between gain and loss frames.

\textbf{Framework Translation.}
\begin{equation}
\Psi_{frame} \subset \Psi_K \quad \text{with} \quad \frac{\partial \text{Choice}}{\partial \Psi_{frame}} \neq 0
\end{equation}
Framing is a context variable that operates through information presentation.

\paragraph{Paper 4: Tversky \& Kahneman (1991)}

``Loss Aversion in Riskless Choice: A Reference-Dependent Model.''
\textit{Quarterly Journal of Economics}, 106(4), 1039--1061.

\textbf{Summary.}
Extends loss aversion beyond risky choice to riskless contexts. Introduces 
formal model of reference-dependent preferences explaining endowment effect, 
status quo bias, and asymmetric price elasticities.

\textbf{Framework Translation.}
Loss aversion is general, not domain-specific:
\begin{equation}
\lambda_d > 1 \quad \forall d \in \{F, E, P, S, D, E\}
\end{equation}

\paragraph{Paper 5: Tversky \& Kahneman (1992)}

``Advances in Prospect Theory: Cumulative Representation of Uncertainty.''
\textit{Journal of Risk and Uncertainty}, 5(4), 297--323.

\textbf{Summary.}
Cumulative Prospect Theory (CPT) extends original PT to handle multiple outcomes.
Introduces rank-dependent probability weighting resolving violations of 
stochastic dominance.

\textbf{Framework Translation.}
Provides the mature version of PT utility:
\begin{equation}
U_{CPT} = \sum_{i: x_i \geq r} \pi^+(p_i) v(x_i - r) + \sum_{i: x_i < r} \pi^-(p_i) v(x_i - r)
\end{equation}

\subsubsection{Loss Aversion and Endowment Effects}

\paragraph{Paper 6: Kahneman, Knetsch \& Thaler (1990)}

``Experimental Tests of the Endowment Effect and the Coase Theorem.''
\textit{Journal of Political Economy}, 98(6), 1325--1348.

\textbf{Summary.}
Classic mug experiments demonstrating WTA/WTP disparity of approximately 2:1.
Shows endowment effect persists even with market experience and incentives.

\textbf{Framework Translation.}
Provides direct measurement of $\lambda$:
\begin{equation}
\lambda = \frac{WTA}{WTP} \approx 2
\end{equation}

\paragraph{Paper 7: Kahneman, Knetsch \& Thaler (1991)}

``Anomalies: The Endowment Effect, Loss Aversion, and Status Quo Bias.''
\textit{Journal of Economic Perspectives}, 5(1), 193--206.

\textbf{Summary.}
Synthesizes evidence on three related phenomena (endowment effect, loss aversion, 
status quo bias) as manifestations of reference-dependent preferences.

\textbf{Framework Translation.}
$C^*$ is sticky because moving away from reference point triggers loss aversion:
\begin{equation}
\frac{\partial C}{\partial t} \to 0 \quad \text{when change requires perceived losses}
\end{equation}

\subsubsection{Synthesis Works}

\paragraph{Paper 8: Kahneman (2003)}

``Maps of Bounded Rationality: Psychology for Behavioral Economics.''
\textit{American Economic Review}, 93(5), 1449--1475.

\textbf{Summary.}
Nobel Prize lecture synthesizing the research program. Introduces accessibility 
(what comes to mind) as unifying concept. Connects heuristics, biases, and 
prospect theory.

\textbf{Framework Translation.}
System 1 accessibility determines which $C_{ij}$ are salient:
\begin{equation}
C_{ij}^{perceived} = f(C_{ij}^{actual}, \text{Accessibility}(\Psi_C, \Psi_K))
\end{equation}

\paragraph{Paper 9: Kahneman (2011)}

\textit{Thinking, Fast and Slow.}
New York: Farrar, Straus and Giroux.

\textbf{Summary.}
Comprehensive synthesis of dual-process theory. System 1 (intuitive) and 
System 2 (deliberative) explain when behavior approaches or deviates from 
rational benchmarks.

\textbf{Framework Translation.}
\begin{equation}
K = K^* \cdot \mathbb{1}_{System2} + K_{biased} \cdot \mathbb{1}_{System1}
\end{equation}
Coherence depends on which system dominates—itself a function of $\Psi_C$.

\paragraph{Paper 10: Kahneman \& Tversky, eds. (2000)}

\textit{Choices, Values, and Frames.}
Cambridge University Press.

\textbf{Summary.}
Definitive edited volume collecting the major papers with commentary. 
Includes extensions by Thaler, Loewenstein, Rabin, and others.

\textbf{Framework Translation.}
Establishes behavioral economics as coherent paradigm with:
\begin{itemize}
    \item Descriptive accuracy (vs. normative EU theory)
    \item Quantified parameters ($\lambda$, $\alpha$, $\beta$)
    \item Policy relevance (choice architecture)
\end{itemize}

%%%%%%%%%%%%%%%%%%%%%%%%%%%%%%%%%%%%%%%%%%%%%%%%%%%%%%%%%%%%%%%%%%%%%%%%%%%%%%%
\subsection{Citation and Verification Note}

Daniel Kahneman (1934--2024): Nobel Prize in Economics 2002, Presidential Medal 
of Freedom 2013. Google Scholar citations $\sim$850,000.

Amos Tversky (1937--1996): Would have shared Nobel Prize (awarded posthumously 
impossible). Grawemeyer Award 1991.

``Prospect Theory'' (1979): Most cited paper in \textit{Econometrica} history.
Most cited economics paper overall per Google Scholar.

%%%%%%%%%%%%%%%%%%%%%%%%%%%%%%%%%%%%%%%%%%%%%%%%%%%%%%%%%%%%%%%%%%%%%%%%%%%%%%%
